\subsection*{Logbook Entry 25: SIM2-centered 21q follow up}

\textbf{What was being measured.}
A follow up was performed on the D3 lower-in-D3 chromosome 21q methylation signal to determine whether the arm-level signal was in fact concentrated into a specific local gene neighborhood. The main aims were to measure unique probe support by annotated gene, quantify genomic compactness of the dominant 21q locus, and generate a first downstream-target scaffold for SIM2-centered interpretation.

\textbf{Methods.}
The input table was
\texttt{results/gse240704/chr21\_locus\_followup/D3\_21q\_probe\_table.csv}.
Unique probe-to-gene membership was reconstructed by splitting manifest-derived gene token strings and counting distinct probes per gene rather than raw repeated token appearances. A SIM2-only subset was then extracted. Genomic compactness was quantified using the observed probe positions on chromosome 21q, including minimum and maximum coordinates, total genomic span, median position, and adjacent probe spacing. Context summaries were retained for refgene-group and CpG-island relation. A seed table of candidate downstream targets of SIM2 was exported as a hypothesis scaffold for later mechanistic overlay work, not as a direct regulatory claim established by this methylation dataset.

\textbf{Results.}
The D3 lower-in-D3 21q signal remained sharply localized. The 21q subset contained 11 probes in total. Unique probe-level gene support showed SIM2 as the dominant annotated locus, with 10 of 11 probes carrying SIM2 annotation, while one probe remained unannotated. The SIM2-centered subset spanned positions 38076869 to 38082613 on chromosome 21, corresponding to a total span of 5744 bp. This indicates that the observed 21q depletion is not diffuse across the arm, but concentrated into a compact local neighborhood around SIM2. The context remained strongly gene-body and CpG-island centered. ERG and RUNX1 were not supported by direct probe annotation in the earlier locus-window follow up, and the present compactness analysis strengthens the interpretation that the 21q signal is best treated as a SIM2-centered local methylation signature rather than an ERG- or RUNX1-centered program.

\textbf{Tables written.}
\texttt{results/gse240704/sim2_followup/D3_21q_probe_gene_membership.csv}\\
\texttt{results/gse240704/sim2_followup/D3_21q_unique_probe_gene_summary.csv}\\
\texttt{results/gse240704/sim2_followup/D3_21q_SIM2_probe_table.csv}\\
\texttt{results/gse240704/sim2_followup/D3_21q_SIM2_compactness_summary.csv}\\
\texttt{results/gse240704/sim2_followup/SIM2_downstream_target_seed_table.csv}\\
\texttt{results/gse240704/sim2_followup/D3_21q_SIM2_followup_summary.json}

\textbf{Figures.}
\texttt{results/gse240704/sim2_followup/plots/D3_21q_SIM2_probe_positions.png}\\
\texttt{results/gse240704/sim2_followup/plots/D3_21q_SIM2_ranked_abs_effect.png}\\
\texttt{results/gse240704/sim2_followup/plots/D3_21q_all_probe_positions.png}

\textbf{Interpretation.}
The chromosome-arm level 21q enrichment can now be refined. In this dataset branch, the lower-in-D3 21q depletion is dominated by a compact SIM2-centered neighborhood, mostly in gene-body and CpG-island contexts. This does not by itself establish SIM2 as the causal driver of the D3 state. However, it does justify replacing a broad arm-only interpretation with a more local 21q neighborhood interpretation centered on SIM2-associated methylation structure.

\textbf{Caveats.}
This step remains annotation dependent. The downstream-target table is only a curated seed scaffold and not a direct target set validated in GSE240704. The analysis is methylation based and does not yet establish transcriptional direction, perturbational dependence, or regulatory causality.

\textbf{Next steps.}
The next step is to connect the SIM2-centered methylation signature to broader D3 biology by testing whether SIM2-related developmental or lineage programs overlap with the existing S-axis biology, especially given the earlier appearance of PTCH1 among positive S-axis loading genes. A later mechanistic overlay step should examine whether SIM2-centered structure aligns with developmental, neural-lineage, or stemness-related modules in a dataset with compatible expression or multi-omic readout.
